The expansion of higher education must be accompanied by questions regarding the application of education in the workplace. Reviews by McGuinness et al. and Somers et al. distinguish between mismatches by level, field and skills: having a job does not prove that the work utilises knowledge related to one's degree \cite{McGuinness2018,Somers2019}. Robst identifies the relationship between one's degree programme and work activities as a specific dimension of labour market integration \cite{Robst2007}. This distinction prevents equating employment outside one's field of study with unemployment or a lack of skills.

The international scale of this issue justifies studying this phenomenon using explicit definitions. The OECD report on the 2023 Survey of Adult Skills, published in 2024, states that, on average across the participating countries and economies, the field of the highest qualification held by 38 per cent of employed people aged 25 to 65 (excluding the self-employed) does not match the field typically required in their job \cite{OECD2024}. This average does not apply to Ecuador or to STEM graduates. It serves to contextualise a problem whose prevalence varies according to the population, classification and method used.

In Ecuador, the 2025 annual ENEMDU survey puts suitable employment at 37.1 per cent of the economically active population and underemployment at 19.4 per cent \cite{InstitutoNacionaldeEstadsticayCensos2026}. These are indicators of the labour market as a whole and not mismatch rates by sector. It is precisely this broad scope that demonstrates why we need a different indicator: even a person in suitable employment may be working in an occupation unrelated to their qualifications. This project will keep income, hours worked, employment status and educational background separate.

The OECD first-cycle PIAAC report places Ecuador approximately 17 percentage points above the OECD average of 40 per cent for field-of-study mismatch \cite{OECD2019}. This implies approximately 57 per cent for the population covered by that indicator, not for STEM graduates specifically. In the annual ENEMDU series, adequate employment rose from 32.5 per cent in 2021 to 37.1 per cent in 2025, while underemployment fell from 23.2 to 19.4 per cent \cite{InstitutoNacionaldeEstadsticayCensos2026}. These trends provide labour-market context, not direct evidence that occupational alignment improved. The gap between general employment indicators and field-specific alignment motivates a separate measure for the STEM population.

The issue has both national and regional roots. Botello Peñaloza examines the educational imbalance in Ecuador; De La Torre and Del Pozo analyse education, skills and income using the ENEMDU; Sierra Pérez and Villarraga Orjuela examine young people in the Andean Community \cite{BotelloPealoza2016,DeLaTorreB2025,SierraPrez2024}. The PIAAC data includes Ecuador and offers another perspective on the utilisation of education, albeit with a different survey and population \cite{OECD2019,OECD2023}. Therefore, the aim of this contribution is not to claim that the country lacks research, but rather to define a predictive classification for employed STEM graduates.

Interest in STEM does not imply a shortage of professionals or homogeneity across degree programmes. Deming and Noray demonstrate the relevance of skills shifts for STEM career paths, whilst Prada and Rucci situate skills development within the context of productive transformations in Latin America \cite{Deming2020,PradaRucci2025}. These studies justify investigating how education and employment are related, but they do not provide grounds for interpreting the number of people in employment as unmet demand or the number of job vacancies. 

The overall objective is to develop and compare machine learning models to classify the mismatch by field of study amongst employed university graduates with STEM degrees in Ecuador during 2021--2025. Specifically, the aim is to standardise the ENEMDU microdata and construct a documented indicator; to compare logistic regression, Random Forest and XGBoost using temporal evaluation and cross-validation; and to identify socioeconomic, demographic and educational variables associated with the classification, examining stability across years and measurement rules.

\noindent\textbf{Research Question (RQ):} How accurately can these models classify the mismatch in this population, and which variables contribute most to that classification?